\documentclass[10pt]{article}
\usepackage[a4paper,margin=2cm]{geometry}
\usepackage{amsmath,amssymb,amsthm}
\usepackage{mathpartir}
\usepackage{xcolor}
\usepackage{booktabs}
\usepackage{hyperref}

\newcommand{\pre}[1]{{\color{gray}#1}}  % presupposition premises (not in the paper)
\newcommand{\U}{\mathsf{U}}
\newcommand{\lam}{\lambda}
\newcommand{\app}{\mathsf{app}}
\newcommand{\strip}{\mathsf{strip}}
\newcommand{\ty}{\ \mathsf{type}}
\newcommand{\TR}{T_R}
\newcommand{\TP}{T_P}
\newcommand{\WN}{\mathsf{WN}_P}
\newcommand{\agda}[1]{\texttt{#1}}

\newtheorem{theorem}{Theorem}
\newtheorem{lemma}[theorem]{Lemma}
\theoremstyle{remark}
\newtheorem{remark}[theorem]{Remark}

\title{Section 4.3 of Sterbac's paper in Agda:\\
annotated versus unannotated Russell universes}
\author{Formalisation notes (\texttt{\textasciitilde/DOMAIN/ERT})}
\date{24--25 September 2026}

\begin{document}
\sloppy
\maketitle

\noindent
This note records the systems $\TR$ (Russell universes, annotated $\lam$ and
$\app$) and $\TP$ (unannotated, ``pure type system'') as formalised in Agda,
how they differ from \texttt{sterbac1.pdf} (§4.3, appendices A.3/A.4; the same
material is §6 of \texttt{sterbac3.pdf}), what is proved, and a gap in the
paper's proof of Lemma~4.17.  Everything is in Agda with
\verb|--without-K --exact-split|, without postulates and without
termination pragmas.  Lemma~4.18 and Theorem~4.19 are proved
\emph{without normalisation} (\S\ref{sec:nf}), from $\Pi$-injectivity,
injectivity of universes and $\Pi/\U$ no-confusion only---the three
properties the paper itself assumes (Theorems 7--9 of \texttt{sterbac3.pdf}).
An earlier proof under weak $\beta$-normalisation of $\TP$ (\S\ref{sec:wn}) is
kept for comparison.  The type-in-type version, where normalisation fails,
is treated in \S\ref{sec:uu}.

Premises printed in \pre{grey} are \emph{presuppositions}: they are not in the
paper's rules, they are admissible (derivable from the other premises), and
they are there only to make the metatheory, the adequacy proof and the
section of $\strip$ structurally recursive (see \S\ref{sec:diff}).

%======================================================================
\section{Syntax}

\[
\begin{array}{lrcl}
\TR: & A,B,a,b,c & ::= & v_i \mid \U_l \mid \Pi(A,B) \mid \lam(A,B,b) \mid \app(A,B,c,a)\\[2pt]
\TP: & A,B,a,b,c & ::= & v_i \mid \U_l \mid \Pi\,A\,B \mid \lam A\, b \mid c\,a
\end{array}
\]
de Bruijn indices, contexts $\Gamma ::= () \mid \Gamma.A$.  The syntax of
$\TP$ is the core syntax of the finite-element model
(\agda{Model/Core.agda}).  The forgetful map is
\[
\strip(v_i)=v_i,\quad \strip(\U_l)=\U_l,\quad
\strip(\Pi(A,B))=\Pi\,\strip(A)\,\strip(B),
\]
\[
\strip(\lam(A,B,b))=\lam\,\strip(A)\,\strip(b),\qquad
\strip(\app(A,B,c,a))=\strip(c)\,\strip(a).
\]
It forgets the codomain annotation of $\lam$ and both annotations of $\app$
(this is Sterbac's $\strip$; it keeps less than Streicher's $\mathsf{strip}_1$,
which keeps $\lam$'s annotation, and than the erasure $|\cdot|$ of §4.2, which
keeps all annotations of $\app$).  $\strip$ commutes with renaming,
substitution and context lookup.

%======================================================================
\section{The system $\TR$ (\agda{RussellTyping.agda})}

Judgements: $\vdash\Gamma$, \ $\Gamma\vdash A\ty$, \ $\Gamma\vdash a:A$, \
$\Gamma\vdash A=B$, \ $\Gamma\vdash a=b:A$.

\paragraph{Contexts and types.}
\begin{mathpar}
\inferrule{ }{\vdash ()} \and
\inferrule{\Gamma\vdash A\ty}{\vdash \Gamma.A} \and
\inferrule{\Gamma\vdash A:\U_l}{\Gamma\vdash A\ty}
\end{mathpar}

\paragraph{Typing.}
\begin{mathpar}
\inferrule{\vdash\Gamma}{\Gamma\vdash v_i:\Gamma(i)} \and
\inferrule{\Gamma\vdash a:A \\ \Gamma\vdash A=B}{\Gamma\vdash a:B} \and
\inferrule{\vdash\Gamma}{\Gamma\vdash \U_l:\U_{l+1}} \and
\inferrule{\Gamma\vdash A:\U_l}{\Gamma\vdash A:\U_{l+1}} \\
\inferrule{\Gamma\vdash A:\U_l \\ \Gamma.A\vdash B:\U_l}{\Gamma\vdash\Pi(A,B):\U_l} \and
\inferrule{\pre{\Gamma\vdash A\ty} \\ \pre{\Gamma.A\vdash B\ty} \\ \Gamma.A\vdash b:B}
          {\Gamma\vdash\lam(A,B,b):\Pi(A,B)} \\
\inferrule{\pre{\Gamma\vdash A\ty} \\ \pre{\Gamma.A\vdash B\ty} \\ \Gamma\vdash c:\Pi(A,B) \\ \Gamma\vdash a:A}
          {\Gamma\vdash\app(A,B,c,a):B[a]}
\end{mathpar}

\paragraph{Type equality.}
\begin{mathpar}
\inferrule{\Gamma\vdash A\ty}{\Gamma\vdash A=A} \and
\inferrule{\Gamma\vdash A=B}{\Gamma\vdash B=A} \and
\inferrule{\Gamma\vdash A=B \\ \Gamma\vdash B=C}{\Gamma\vdash A=C} \and
\inferrule{\Gamma\vdash A=B:\U_l}{\Gamma\vdash A=B} \\
\inferrule{\pre{\Gamma\vdash A\ty} \\ \pre{\Gamma.A\vdash B\ty} \\ \Gamma\vdash A=A' \\ \Gamma.A\vdash B=B'}
          {\Gamma\vdash\Pi(A,B)=\Pi(A',B')}
\end{mathpar}

\paragraph{Term equality.}
\begin{mathpar}
\inferrule{\Gamma\vdash a:A}{\Gamma\vdash a=a:A} \and
\inferrule{\Gamma\vdash a=b:A}{\Gamma\vdash b=a:A} \and
\inferrule{\Gamma\vdash a=b:A \\ \Gamma\vdash b=c:A}{\Gamma\vdash a=c:A} \and
\inferrule{\Gamma\vdash a=b:A \\ \Gamma\vdash A=B}{\Gamma\vdash a=b:B} \and
\inferrule{\Gamma\vdash A=B:\U_l}{\Gamma\vdash A=B:\U_{l+1}} \\
\inferrule{\pre{\Gamma\vdash A:\U_l} \\ \pre{\Gamma.A\vdash B:\U_l} \\ \Gamma\vdash A=A':\U_l \\ \Gamma.A\vdash B=B':\U_l}
          {\Gamma\vdash\Pi(A,B)=\Pi(A',B'):\U_l} \\
\inferrule{\pre{\Gamma\vdash A\ty} \\ \pre{\Gamma.A\vdash B\ty} \\ \pre{\Gamma.A\vdash b:B} \\ \Gamma.A\vdash b=b':B}
          {\Gamma\vdash\lam(A,B,b)=\lam(A,B,b'):\Pi(A,B)} \\
\inferrule{\pre{\Gamma\vdash A\ty} \\ \pre{\Gamma.A\vdash B\ty} \\ \Gamma\vdash A=A' \\ \Gamma.A\vdash B=B' \\ \Gamma.A\vdash b:B}
          {\Gamma\vdash\lam(A,B,b)=\lam(A',B',b):\Pi(A,B)} \\
\inferrule{\pre{\Gamma\vdash A\ty} \\ \pre{\Gamma.A\vdash B\ty} \\ \Gamma\vdash c=c':\Pi(A,B) \\ \Gamma\vdash a:A}
          {\Gamma\vdash\app(A,B,c,a)=\app(A,B,c',a):B[a]} \\
\inferrule{\pre{\Gamma\vdash A\ty} \\ \pre{\Gamma.A\vdash B\ty} \\ \Gamma\vdash c:\Pi(A,B) \\ \Gamma\vdash a=a':A \\ \pre{\Gamma\vdash B[a]=B[a']}}
          {\Gamma\vdash\app(A,B,c,a)=\app(A,B,c,a'):B[a]} \\
\inferrule{\pre{\Gamma\vdash A\ty} \\ \pre{\Gamma.A\vdash B\ty} \\ \Gamma\vdash A=A' \\ \Gamma.A\vdash B=B' \\ \Gamma\vdash c:\Pi(A,B) \\ \Gamma\vdash a:A}
          {\Gamma\vdash\app(A,B,c,a)=\app(A',B',c,a):B[a]} \\
\inferrule{\pre{\Gamma\vdash A\ty} \\ \pre{\Gamma.A\vdash B\ty} \\ \Gamma.A\vdash b:B \\ \Gamma\vdash a:A}
          {\Gamma\vdash\app(A,B,\lam(A,B,b),a)=b[a]:B[a]}\ (\beta) \and
\inferrule{\pre{\Gamma\vdash A\ty} \\ \pre{\Gamma.A\vdash B\ty} \\ \Gamma\vdash c:\Pi(A,B)}
          {\Gamma\vdash c=\lam(A,B,\app(A\!\uparrow,B\!\uparrow^{+},c\!\uparrow,v_0)):\Pi(A,B)}\ (\eta)
\end{mathpar}

%======================================================================
\section{The system $\TP$ (\agda{PTyping.agda})}

The rules of $\TP$ are those of $\TR$ with the annotations removed:
contexts, types, type equality and the structural rules are as above
(on $\TP$ syntax), and
\begin{mathpar}
\inferrule{\Gamma\vdash A:\U_l \\ \Gamma.A\vdash B:\U_l}{\Gamma\vdash\Pi\,A\,B:\U_l} \and
\inferrule{\pre{\Gamma\vdash A\ty} \\ \pre{\Gamma.A\vdash B\ty} \\ \Gamma.A\vdash b:B}{\Gamma\vdash\lam A\,b:\Pi\,A\,B} \and
\inferrule{\pre{\Gamma\vdash A\ty} \\ \pre{\Gamma.A\vdash B\ty} \\ \Gamma\vdash c:\Pi\,A\,B \\ \Gamma\vdash a:A}{\Gamma\vdash c\,a:B[a]} \\
\inferrule{\pre{\Gamma\vdash A\ty} \\ \pre{\Gamma.A\vdash B\ty} \\ \Gamma\vdash A=A' \\ \Gamma.A\vdash B=B'}{\Gamma\vdash\Pi\,A\,B=\Pi\,A'\,B'} \and
\inferrule{\pre{\Gamma\vdash A:\U_l} \\ \pre{\Gamma.A\vdash B:\U_l} \\ \Gamma\vdash A=A':\U_l \\ \Gamma.A\vdash B=B':\U_l}{\Gamma\vdash\Pi\,A\,B=\Pi\,A'\,B':\U_l} \\
\inferrule{\pre{\Gamma\vdash A\ty} \\ \pre{\Gamma.A\vdash B\ty} \\ \pre{\Gamma.A\vdash b:B} \\ \Gamma.A\vdash b=b':B}{\Gamma\vdash\lam A\,b=\lam A\,b':\Pi\,A\,B} \and
\inferrule{\pre{\Gamma\vdash A\ty} \\ \pre{\Gamma.A\vdash B\ty} \\ \Gamma\vdash A=A' \\ \Gamma.A\vdash b:B}{\Gamma\vdash\lam A\,b=\lam A'\,b:\Pi\,A\,B} \\
\inferrule{\pre{\Gamma\vdash A\ty} \\ \pre{\Gamma.A\vdash B\ty} \\ \Gamma\vdash c=c':\Pi\,A\,B \\ \Gamma\vdash a:A}{\Gamma\vdash c\,a=c'\,a:B[a]} \and
\inferrule{\pre{\Gamma\vdash A\ty} \\ \pre{\Gamma.A\vdash B\ty} \\ \Gamma\vdash c:\Pi\,A\,B \\ \Gamma\vdash a=a':A}{\Gamma\vdash c\,a=c\,a':B[a]} \\
\inferrule{\pre{\Gamma\vdash A\ty} \\ \pre{\Gamma.A\vdash B\ty} \\ \Gamma.A\vdash b:B \\ \Gamma\vdash a:A}{\Gamma\vdash (\lam A\,b)\,a=b[a]:B[a]}\ (\beta) \and
\inferrule{\pre{\Gamma\vdash A\ty} \\ \pre{\Gamma.A\vdash B\ty} \\ \Gamma\vdash c:\Pi\,A\,B}{\Gamma\vdash c=\lam A\,(c\!\uparrow v_0):\Pi\,A\,B}\ (\eta)
\end{mathpar}
There is no counterpart of the $\TR$ rule changing the annotations of an
application: its two sides have the same stripping.

%======================================================================
\section{Differences with the paper}\label{sec:diff}

\begin{enumerate}
\item \textbf{Presupposition premises} (grey).  The paper's rules have
  none.  In $\TR$ they are needed at two places:
  \begin{itemize}
  \item the adequacy proof of the domain model (for $\Pi$-injectivity) is a
    structural induction whose congruence cases need the
    \emph{two-substitution} validity of the codomain / body, which must come
    from a subderivation (this is the same obstruction as MIN's
    \texttt{conv-Pi});
  \item Theorem~4.15 and Theorem~4.19 are structural inductions which need
    the typings of the components of binders (to lift the extended context).
  \end{itemize}
  The paper's rules are admissible in $\TR$: they are the lemmas
  \agda{mk-conv-Ty-Pi}, \agda{mk-conv-cong-Pi}, \agda{mk-conv-cong-Lam-Ty},
  \agda{mk-conv-cong-App-Ty} (\agda{RussellMeta.agda}), using the
  presupposition lemmas.  For $\TP$ the corresponding admissibility is not
  formalised (it needs presupposition lemmas for $\TP$).
\item \textbf{Universes and cumulativity} are one step at a time
  ($\U_l:\U_{l+1}$, $A:\U_l\Rightarrow A:\U_{l+1}$); the paper's
  $\U_m:\U_n$ ($m<n$) and $A:\U_n\Rightarrow A:\U_p$ ($n<p$) are derivable
  (\agda{cum-le}).  $\Gamma\vdash\U_n\ty$ is derived.
\item \textbf{Variables} are typed by lookup ($\Gamma\vdash v_i:\Gamma(i)$
  with the weakened type) instead of the paper's two weakening rules.
\item The congruences of $\lam$ and $\app$ are split into one rule per
  argument (body / annotation, function / argument / annotations).
\end{enumerate}

%======================================================================
\section{What is proved without any hypothesis}

\begin{theorem}[Adequacy of the domain model, \agda{Adeq/Driver.agda}]
The finite-element (Coquand--Huber) model with levelled universe codes is
adequate for $\TR$: every derivation is valid for the rank-stratified
logical relation, by structural induction.
\end{theorem}

\begin{theorem}[\agda{PiInjectivityR.agda}, \agda{RussellInversion.agda}]
In $\TR$: if $\Gamma\vdash\Pi(C_0,B_0)=\Pi(C_1,B_1)$ then
$\Gamma\vdash C_0=C_1$ and $\Gamma.C_0\vdash B_0=B_1$ ($\Pi$-injectivity,
from adequacy); if $\Gamma\vdash\U_m=\U_k$ then $m=k$; and
$\Gamma\vdash\Pi(A,B)=\U_m$ never holds (from soundness of the model).
\end{theorem}

\begin{theorem}[4.15, \agda{StripDeriv.agda}]
If $\Gamma\vdash J$ in $\TR$ then $\strip(\Gamma)\vdash\strip(J)$ in $\TP$.
\end{theorem}

\begin{theorem}[Subject reduction, \agda{SubjectReductionR.agda}]
Let $u\to u'$ be a $\beta$-step of $\TR$ at a position kept by $\strip$
($\app(A,B,\lam(A',B',b),a)\to b[a]$, annotations arbitrary).  If
$\Gamma\vdash u:T$ then $\Gamma\vdash u=u':T$.  (The $\beta$ case uses
$\Pi$-injectivity.)
\end{theorem}

\begin{lemma}[Lifting of reductions, \agda{RussellStep.agda}]
If $\strip(u)\to M'$ is a $\beta$-step of $\TP$, then $u\to u'$ in $\TR$ with
$\strip(u')=M'$, for \emph{every} $u$.
\end{lemma}

\begin{lemma}[Normal forms, \agda{StripUniq.agda}: \agda{nf-conv}]
Let $\Gamma\vdash u_0:A_0$, $\Gamma\vdash u_1:A_1$ in $\TR$ with
$\strip(u_0)=\strip(u_1)$ a $\beta$-normal form, and $\Gamma\vdash A_0=A_1$.
Then $\Gamma\vdash u_0=u_1:A_0$.  Similarly two types with the same
$\beta$-normal stripping are equal.
\end{lemma}
\noindent
The proof is by induction on the normal form, together with: for
\emph{neutral} strippings, $u_0,u_1$ are equal at a common type $P$ lying below
both $A_0$ and $A_1$, where ``$P$ below $T$'' means $\Gamma\vdash P=T$, or
$\Gamma\vdash P=\U_k$ and $\Gamma\vdash T=\U_m$ with $k\le m$
(\agda{SubTy}).  In a neutral application the annotations are forced by the
type of the head; in a $\lam$ by its type (via $\Pi$-injectivity);
cumulativity only acts at the root, where it changes no syntax.


%======================================================================
\section{Lemma 4.18 without normalisation}\label{sec:nf}

\agda{StripUniqNF.agda}; the only non-syntactic ingredients are
$\Pi$-injectivity, injectivity of universes and $\Pi/\U$ no-confusion.
The rest is a structural induction on the term, i.e.\ a primitive recursive
transformation of derivations.

\paragraph{What the induction hypothesis must say.}  In the counterexample
of \S\ref{sec:gap}, $f_0=\lam(\U_1,\U_1,v_0)$ and $f_1=\lam(\U_1,\U_2,v_0)$
have the same stripping but are \emph{not} convertible: their types
$\Pi(\U_1,\U_1)$, $\Pi(\U_1,\U_2)$ differ.  So the hypothesis has to relate
terms whose types differ by cumulativity \emph{under} $\Pi$.  $\TR$ has no
such subtyping, but it is definable by $\eta$-expansion:
\[
\mathsf{coe}_{\Pi(A,B)\le\Pi(A,L)}(t) \;=\;
\lam\bigl(A,\,L,\,\mathsf{coe}_{B\le L}(\app(A{\uparrow},B{\uparrow^+},t{\uparrow},v_0))\bigr),
\qquad \mathsf{coe}_{\U_k\le\U_m}(t)=t .
\]
\paragraph{Joins.}  $\mathsf{Join}_\Gamma(T_0,T_1,L)$ is generated by
\begin{itemize}
\item $\Gamma\vdash T_0=L$ and $\Gamma\vdash T_1=L$ (coercions: identity);
\item $\Gamma\vdash T_0=\U_k$, $\Gamma\vdash T_1=\U_m$, $L=\U_{\max(k,m)}$
  (identity; this is where cumulativity enters);
\item $\Gamma\vdash T_i=\Pi(A,B_i)$ and $\mathsf{Join}_{\Gamma.A}(B_0,B_1,L')$,
  $L=\Pi(A,L')$ (coercions as above).
\end{itemize}
A join carries its coercions $\mathsf{coe}_0$, $\mathsf{coe}_1$.  They are
well typed ($t:T_i\Rightarrow\mathsf{coe}_i(t):L$), congruent, stable
under substitution (syntactically: $\mathsf{coe}_k(t)[\sigma]=
\mathsf{coe}_{k[\sigma]}(t[\sigma])$), and compute under application:
$\app(A,L,\mathsf{coe}_{\Pi(A,B)\le\Pi(A,L)}(t),a)=\mathsf{coe}_{B[a]\le L[a]}(\app(A,B,t,a))$ by $\beta$.

\begin{lemma}[main]
If $\Gamma\vdash u_0:T_0$, $\Gamma\vdash u_1:T_1$ and
$\strip(u_0)=\strip(u_1)$, then there is a join
$\mathsf{Join}_\Gamma(T_0,T_1,L)$ with
$\Gamma\vdash\mathsf{coe}_0(u_0)=\mathsf{coe}_1(u_1):L$.
\end{lemma}
\noindent
The proof is by induction on $u_0$, with $T_0,T_1$ \emph{arbitrary} derivable
types.  Inversion gives the principal type $P_i$ of $u_i$ with
$P_i\le T_i$ (conversion, or universes $\U_k\le\U_m$).  A join of $P_0,P_1$
is then moved to one of $T_0,T_1$: by transport along conversions, or,
if a universe is involved, by a universe join of larger level.
\begin{itemize}
\item $v_i$, $\U_l$: trivial.  $\Pi(A_i,B_i)$: the induction hypothesis on
  the components gives joins of universes, hence $A_0=A_1:\U_m$ and
  $B_0=B_1:\U_m$.
\item $\lam(A_i,B_i,b_i)$: $A_0=A_1$ as above; the induction hypothesis on
  the bodies, at their annotated types $B_0,B_1$, gives
  $\mathsf{Join}_{\Gamma.A_0}(B_0,B_1,L')$.  The $\Pi$-join
  $\mathsf{Join}(\Pi(A_0,B_0),\Pi(A_1,B_1),\Pi(A_0,L'))$ works because
  $\mathsf{coe}(\lam(A,B_0,b_0))=\lam(A,L',\mathsf{coe}(b_0))$ by $\beta$
  under the binder.
\item $\app(C_i,D_i,c_i,a_i)$: the induction hypothesis on $c_0,c_1$ at
  $\Pi(C_i,D_i)$ gives a join which, by no-confusion and $\Pi$-injectivity,
  is either a conversion or a $\Pi$-join $\Pi(A,L')$ with $C_i=A$, $D_i=B_i$.
  The induction hypothesis on $a_0,a_1$ together with the diagonal lemma
  below gives $a_0=a_1:C_0$.  Applying the coercion equation to $a_0$ and
  using the $\beta$-law of coercions gives
  $\mathsf{coe}_{B_0[a_0]\le L'[a_0]}(u_0)=\mathsf{coe}_{B_1[a_0]\le L'[a_0]}(u_1):L'[a_0]$,
  with the substituted join $\mathsf{Join}(B_0[a_0],B_1[a_0],L'[a_0])$.
  \emph{This is the step the paper's proof misses}: the functions are
  compared at the join of their types, not at a common type.
\end{itemize}

\begin{lemma}[diagonal]
If $\mathsf{Join}_\Gamma(T_0,T_1,L)$, $\Gamma\vdash T_0=T_1$, $t_i:T_i$ and
$\mathsf{coe}_0(t_0)=\mathsf{coe}_1(t_1):L$, then $\Gamma\vdash t_0=t_1:T_0$.
\end{lemma}
\noindent
By induction on the join.  In the universe case $k=m$ by injectivity of
universes, so $L=\U_{\max(k,k)}=\U_k$.  No lowering of a conversion from a
larger universe is ever needed: this is why the level of a universe join is
exactly $\max(k,m)$.  In the $\Pi$ case, $B_0=B_1$ by $\Pi$-injectivity.
Applying both $\lam$'s to $v_0$ (after weakening) and using $\beta$, the
coerced bodies are equal; then the induction hypothesis applies, and $\eta$
concludes.

\begin{theorem}[4.18, no normalisation]
If $\Gamma\vdash u_0:A_0$, $\Gamma\vdash u_1:A_1$, $\strip(u_0)=\strip(u_1)$ and
$\Gamma\vdash A_0=A_1$, then $\Gamma\vdash u_0=u_1:A_0$
(\agda{term-uniq-conv-NF}); similarly for types (\agda{type-uniq-NF}), and
Lemma~4.18 (\agda{term-uniq-NF}).  Consequently Theorem~4.19
(\agda{SectionR.agda}) holds without normalisation as well.
\end{theorem}
\noindent
\agda{StripUniqNFTest.agda} checks the counterexample of \S\ref{sec:gap}:
$\vdash u_0=u_1:\U_2$ follows from the theorem.

\begin{remark}
$\eta$ is used essentially: it is what makes coercions injective.
The argument is a structural induction producing derivations from
derivations; it uses no logical relation and no normalisation.  The three
properties it takes as input are proved here from the domain model; the
paper assumes them (from normalisation of the GAT, or from the finite-domain
model).
\end{remark}


%======================================================================
\section{The case $\U:\U$ (\texttt{ERTUU/})}\label{sec:uu}

\texttt{ERTUU/} is the type-in-type version: one universe, $\U:\U$,
no levels and no cumulativity.  It contains §4.2 (Tarski
$\Leftrightarrow\TR$) and now also §4.3 ($\TR\Leftrightarrow\TP$).
Everything is checked with \verb|--safe| (\agda{ERTUU/All.agda},
72 modules, about 13\,s from scratch).

\paragraph{Why this case is interesting.}  $\TP$ with $\U:\U$ does
\emph{not} normalise (Girard's paradox), so the route through
$\beta$-normal forms (\S\ref{sec:wn}) is impossible.  The only
non-syntactic ingredient, $\Pi$-injectivity of $\TR$, comes from the
domain model, and its adequacy proof does not need normalisation.

\begin{theorem}[\agda{ERTUU/PiInjectivityR.agda}]
In $\TR$ with $\U:\U$: if $\Gamma\vdash\Pi(C_0,B_0)=\Pi(C_1,B_1)$ then
$\Gamma\vdash C_0=C_1$ and $\Gamma.C_0\vdash B_0=B_1$.
\end{theorem}
\noindent
The model (\texttt{Dom/}, \texttt{Model/}) is the finite-element model
of \texttt{ERT/}, with levelled universe codes.  $\U$ is interpreted
as the level-0 universe code.  This is sound because the model's
membership relation ignores levels: soundness
(\agda{Model/RussellSound}) goes through with $\U:\U$ interpreted as
``$\U_0:\U_0$''.  The validity relation and the adequacy driver
(\texttt{Valid/}, \texttt{Adeq/}) are those of \texttt{ERT/} with
the universe levels removed from the syntax.  The only change needed by
hand is that $\U$ now evaluates to the level-0 code only
(\agda{Adeq/Stmt.evU}).

\begin{theorem}[4.15, 4.17/4.18, 4.19 for $\U:\U$]
\agda{StripDeriv}: $\TR$ derivations strip to $\TP$ derivations.
\agda{StripUniq}: if $\Gamma\vdash u_0:T_0$, $\Gamma\vdash u_1:T_1$ and
$\strip(u_0)=\strip(u_1)$, then $\Gamma\vdash T_0=T_1$ and
$\Gamma\vdash u_0=u_1:T_0$; no hypothesis on the types is needed.
\agda{SectionR}: every $\TP$ judgement lifts into every $\TR$ context
over its context, uniquely up to conversion.
\end{theorem}
\noindent
Without cumulativity, types are unique up to conversion, and the paper's
proof of Lemma 4.17 (induction on the term, with $\Pi$-injectivity in the
application case) works as written.  The joins and coercions of
\S\ref{sec:nf} are not needed: the counterexample of \S\ref{sec:gap}
relies on cumulativity.  Neither injectivity of universes nor
no-confusion is used.

The $\TR$ rules of \texttt{ERTUU/} now carry the same presupposition
premises as those of \texttt{ERT/} (\S\ref{sec:diff}); the §4.2 files
(\agda{EraseDeriv}, \agda{Equivalence}) use the paper-form lemmas
\agda{mk-*}.  The \texttt{CoverageNoExactSplit} warning in
\agda{ERTUU/Uniqueness.term-uniq} (§4.2) was harmless: a catch-all
clause for a conversion on the right overlapped the left-hand conversion
clause, so it failed to hold definitionally only when both sides were
conversions.  It is now split into one clause per left-hand head, and
the whole of \texttt{ERTUU/} checks with no warnings.  (The analogous
warning in \agda{ERT/UniquenessTermPartial} is of the same benign
kind and is still present.)

%======================================================================
\section{What is proved assuming weak $\beta$-normalisation (superseded)}\label{sec:wn}

\[
\WN:\quad \text{if } \Gamma\vdash M:A \text{ in } \TP \text{ then }
M\to^*_\beta N \text{ for some } \beta\text{-normal } N
\]
($\beta$ at all positions of $\TP$ terms; \agda{PReduction.agda}).  It is a
module parameter, not a postulate.

\begin{theorem}[4.18, \agda{StripUniq.agda}]
Assume $\WN$.
\begin{enumerate}
\item If $\Gamma\vdash A_0\ty$, $\Gamma\vdash A_1\ty$ in $\TR$ and
  $\strip(A_0)=\strip(A_1)$, then $\Gamma\vdash A_0=A_1$
  (\agda{type-uniq-R}).
\item If $\Gamma\vdash u_0:A_0$, $\Gamma\vdash u_1:A_1$ and
  $\strip(u_0)=\strip(u_1)$ and $\Gamma\vdash A_0=A_1$, then
  $\Gamma\vdash u_0=u_1:A_0$ (\agda{term-uniq-conv-R}).  In particular
  (Lemma~4.18) this holds if $\strip(A_0)=\strip(A_1)$ (\agda{term-uniq-R}).
\end{enumerate}
\end{theorem}
\noindent
Proof (Streicher's method, \emph{Semantics of Type Theory}, Thm~4.12):
normalise $\strip(u_0)$, perform the same reduction sequence on $u_0$ and on
$u_1$ (lifting lemma), use subject reduction, and conclude by the normal-form
lemma.  Lemma~4.17 and $\strip'$ are not needed.

\begin{theorem}[4.19, \agda{SectionR.agda}]
Assume $\WN$.  Let $\Gamma\vdash J$ in $\TP$.
\begin{enumerate}
\item For every $\Gamma'$ with $\vdash\Gamma'$ in $\TR$ and
  $\strip(\Gamma')=\Gamma$ there is a derivation $\Gamma'\vdash J'$ in $\TR$
  with $\strip(J')=J$ (\agda{lift-IsType}, \agda{lift-HasType},
  \agda{lift-ConvTy}, \agda{lift-ConvTm}); and such a $\Gamma'$ exists
  (\agda{lift-WfCtx}).
\item Two such $J'$ in the same $\Gamma'$ are equal up to conversion
  (\agda{section-uniq-Tm}, \agda{section-uniq-Ty}, i.e.\ Theorem~4.18).
\end{enumerate}
\end{theorem}
\noindent
The lift is by structural induction on the $\TP$ derivation; annotations
are read off the premises, and premises that share a term or a type are
reconciled by Theorem~4.18.

%======================================================================
\section{The gap in the paper's proof of Lemma 4.17}\label{sec:gap}

\paragraph{The statement} (Lemma 4.17 of \texttt{sterbac1.pdf}, Lemma 30 of
\texttt{sterbac3.pdf}).  With $\strip'(\Pi(A,B))=\strip'(B)$ and $\strip'=\strip$
otherwise: (i) if $\Gamma\vdash A\ty$, $\Gamma\vdash B\ty$ and
$\strip(A)=\strip(B)$ then $\Gamma\vdash A=B$; (ii) if
$\Gamma\vdash u_0:A_0$, $\Gamma\vdash u_1:A_1$, $\strip(u_0)=\strip(u_1)$ and
$\strip'(A_0)=\strip'(A_1)$ then $\Gamma\vdash u_0=u_1:A_0$ and
$\Gamma\vdash A_0=A_1$.  The proof is ``by mutual induction on $A$ and $u_0$''.

\paragraph{The application case} assumes that the last rule of the derivation
of $\Gamma\vdash u_0:A_0$ is the application rule, so that $A_0$ is
$B_0[a_0]$, and derives $\strip'(B_0)=\strip'(B_1)$ from
$\strip'(B_0[a_0])=\strip'(B_1[a_1])$, then applies the induction hypothesis
to the functions to get $\Gamma\vdash\Pi(C_0,B_0)=\Pi(C_1,B_1)$.  In $\TR$
the last rule may be \emph{cumulativity}.  Counterexample
(\agda{Gap417.agda}, checked in Agda), in the empty context:
\[
\begin{array}{ll}
f_0=\lam(\U_1,\U_1,v_0):\Pi(\U_1,\U_1), &
f_1=\lam(\U_1,\U_2,v_0):\Pi(\U_1,\U_2)\quad(v_0:\U_1\le\U_2),\\[2pt]
u_0=\app(\U_1,\U_1,f_0,\U_0):\U_1\le\U_2, &
u_1=\app(\U_1,\U_2,f_1,\U_0):\U_2 .
\end{array}
\]
Here $\strip(u_0)=\strip(u_1)=(\lam\,\U_1\,v_0)\,\U_0$ and
$A_0=A_1=\U_2$, so the hypotheses of (ii) hold.  But
$\strip'(\Pi(\U_1,\U_1))=\U_1\ne\U_2=\strip'(\Pi(\U_1,\U_2))$, so the
induction hypothesis does not apply to $f_0,f_1$, and the conclusion the
proof needs, $\vdash\Pi(\U_1,\U_1)=\Pi(\U_1,\U_2)$, is \emph{false}
($\Pi$-injectivity and injectivity of universes; \agda{no-Pi-conv}).
The statement itself is not refuted: $\vdash u_0=u_1:\U_2$ holds, both sides
$\beta$-reducing to $\U_0$ (\agda{u0=u1}).

\paragraph{Further weak steps} (not formalised as counterexamples).
\begin{itemize}
\item ``$\strip'(B_0[a_0])=\strip'(B_1[a_1])$ implies $\strip'(B_0)=\strip'(B_1)$''
  fails for open codomains, since substitution may identify different
  codomains ($B_0=v_0$, $B_1=\U_1$, $a_0=a_1=\U_1$).
\item In the $\lam$ case the proof uses
  $\strip'(\Pi(C_0,B_0))=\strip(B_0)$, i.e.\ it assumes $A_0$ is literally
  $\Pi(C_0,B_0)$; in $\TR$, $A_0$ is only convertible to it.
\item The induction measure (``on $A$ and $u_0$'') is not made precise, and
  the type part (i) inherits the problem: types may be $\beta$-redexes whose
  $\lam$'s codomain annotations differ by cumulativity.
\end{itemize}

\paragraph{Why Streicher's argument does not transfer.}  In
\emph{Semantics of Type Theory}, Thm~4.11 ($\mathsf{strip}_1$, dropping the
annotations of applications) is proved by induction on the number of
applications, and its application case uses Thm~4.10, i.e.\
\emph{uniqueness of types} (and injectivity of $\Pi$).  $\TR$ has no
uniqueness of types because of cumulativity, and $\strip'$ does not repair
this.  Streicher's Thm~4.12 (dropping also $\lam$'s annotation) proceeds
instead through $\beta$-normal forms; this is the method used above, under
$\WN$.  The paper notices the cumulativity phenomenon only for the further
step of erasing $\lam$'s domain (Example 4.20 / 33), where it states that
normalisation is needed.  Note also that §4.2 (Tarski $\Leftrightarrow$
$\TR$) is not affected: there the erasure keeps the annotations of $\app$,
and the formalised proof (\agda{Uniqueness.agda},
\agda{UniquenessTermPartial.agda}, \agda{Equivalence.agda}) uses no
normalisation.

\paragraph{Consequence.}  The gap is in the proof, not in the method:
Lemma~4.18 and Theorem~4.19 hold with a purely syntactic proof
(\S\ref{sec:nf}), provided the induction hypothesis compares terms at the
\emph{join} of their types.  So the claim of \texttt{sterbac3.pdf} that the
proof is syntactic, relative to injectivity and no-confusion, survives for
§6 too.

%======================================================================
\section*{Agda files}
\begin{tabular}{@{}ll@{}}
\toprule
\agda{RussellTyping}, \agda{RussellMeta}, \agda{RussellMetaCong} & $\TR$ and its metatheory \\
\agda{PTyping}, \agda{PReduction} & $\TP$, $\beta$-reduction, normal forms, $\WN$ \\
\agda{Model/*}, \agda{Dom/*}, \agda{Valid/*}, \agda{Adeq/*} & domain model, soundness, validity, adequacy \\
\agda{PiInjectivityR}, \agda{RussellInversion} & $\Pi$-injectivity, universe injectivity, inversion \\
\agda{StripDeriv} & Theorem 4.15 \\
\agda{RussellStep}, \agda{SubjectReductionR} & lifting of reductions, subject reduction \\
\agda{StripUniqNF} & coercions, joins, Theorem 4.18 (no normalisation) \\
\agda{SectionR} & Theorem 4.19 (no normalisation) \\
\agda{StripUniq} & superseded: Theorem 4.18 under $\WN$ \\
\agda{Gap417} & the counterexample to the proof of Lemma 4.17 \\
\texttt{ERTUU/} & $\U:\U$: §4.2 and §4.3 (\agda{All.agda}, \verb|--safe|) \\
\bottomrule
\end{tabular}

\end{document}
